\section{Experimental Design}

\subsection{General testing methodology}

Each trial uses one unique document and one explicit PyMongo session for all
subject operations. A single logical writer assigns integer versions; update
records retain parent-write and read-dependency identifiers. Setup writes and
diagnostic reads are recorded separately from subject operations. The recorder
captures the requested route, actual server address, directly observed member
role, session and concern settings, command status, operation times, fault
events, and returned document state. Scalar timestamps support diagnosis but
do not replace application-level dependencies in the checkers.

Three actors are separated. The subject client issues the tested history. An
independent topology and document observer connects directly to all members and
does not use the subject session. A restricted fault controller changes only
replica-network reachability. The runner has no Docker socket or Docker
credentials.

\subsection{Consistency property checkers}

The RYW checker compares a later read with the version of an earlier completed
write by the same client. MR compares versions returned by successive reads of
the same key. MW checks that a converged state containing a successor write
also contains its predecessor. WFR checks that a converged state containing a
dependent write includes the version returned by its preceding read. MW and
WFR use the same logical document so that the checker can evaluate the stated
dependency.

After fault cleanup, the topology oracle requires one primary and two
secondaries, and the independent observer checks that all three copies contain
the same complete update list. The checker reports six distinct outcomes:
\texttt{PASS}, \texttt{VIOLATION}, \texttt{UNAVAILABLE},
\texttt{INDETERMINATE}, \texttt{PRECONDITION\_MISS}, and
\texttt{HARNESS\_ERROR}. Only PASS and VIOLATION enter a consistency rate;
timeouts, unresolved effects, missed preconditions, and harness failures remain
separate.

\subsection{Normal-operation scenario}

The normal control runs the same property schedule, configuration, session
rules, and operation deadlines on a stable replica set without an injected
fault. Ten repetitions were run for each of eight configurations and four
properties, for 320 controls. They are reported separately from the 960
adversarial RQ1 histories so that the baseline does not dilute fault-condition
outcomes.

\subsection{Node-failure scenarios}

F1 stops one secondary while the primary and remaining secondary continue to
serve the replica set. Its rationale is to remove redundancy without forcing a
primary election. F2 stops the primary and waits for the two remaining members
to elect a new primary. The trial issues no subject request during this
election barrier, so its scheduled-operation success cannot measure service
availability during the interval. The runner records election and recovery
times by fault episode, then verifies topology and data convergence after
restart.

\subsection{Network-partition scenario}

F3 blocks replication traffic between the former primary and the other two
members while preserving the test client's route to the isolated member. The
connected pair is the majority side. Controller state, direct member access,
observed roles, and version state verify the intended condition; a successful
firewall command alone is not treated as proof of isolation. Target members
rotate across episodes. Healing is followed by a stable-topology check and an
independent three-member convergence check.

\subsection{Experimental procedure and repetitions}

RQ1 ran sequentially with 30 adversarial repetitions for every C1-C8 and
property cell, giving 960 histories, in addition to the 320 normal controls.
The smoke campaign checks schedule construction and is excluded from evidence;
the pilot is also kept separate. RQ2 crosses C1/C3/C4/C6 with F1/F2/F3 and the
four properties for 384 core histories. Four F3 signature cells received 12
extra repetitions each, adding 48 histories. RQ2 therefore contains 432
histories grouped in 36 shared fault episodes. Histories in one episode are
not independent fault injections. RQ2 reuses matching RQ1 normal controls as
its descriptive baseline and does not add or pool another normal campaign.

The property schedules construct the state needed by their checkers:

\begin{description}
\item[RYW] Isolate a stale secondary, confirm it remains at version 0 while a
  fresh replica has version 1, write version 1 through the subject session,
  then read from the stale secondary. The split makes a stale returned value
  observable; a failed read is unavailable, not a consistency violation.
\item[MR] Isolate a stale secondary before setup, create version 1, confirm the
  fresh/stale version split, then read from the fresh and stale secondaries in
  that order through one session. The second read tests for regression.
\item[MW] Isolate the old primary, issue the first subject write there, wait
  for the majority side to elect a new primary, then issue its dependent
  successor write. The post-recovery snapshot checks whether the successor
  remains without its predecessor.
\item[WFR] Write version 1 on the isolated old branch as setup, read that
  concrete version through the subject session, wait for the majority-side
  election, and issue a dependent write on the new primary. The final snapshot
  tests whether the dependent write retained the read version.
\end{description}

These schedules are designed to expose the named property, not to apply one
generic fault sequence to every history. Preconditions are checked before
subject operations. A dependent operation is not issued when its prerequisite
did not complete or return a concrete version.

\subsection{Measurements and outcome classification}

The analysis reports history outcomes, the consistency rate over PASS plus
VIOLATION, subject-operation success and history completion, and operation
latency p50/p95/p99. RQ2 additionally reports acknowledged-write rollback,
election, recovery, and fault-application/convergence counts. Fault timing uses
episodes as its unit; latency includes timed-out operations when start and end
times exist. All deadlines, retry settings, and precondition rules are
registered in \path{docs/experimental-protocol.md}. The representative
operation schedules are shown in Figure~\ref{fig:property-timelines}.
